\section{ESTADO DEL SENSOR INTERNO}
\begin{table}[H]
    \begin{tblr}{
        cell{6-7,10-12,14-15}{1} = {font=\footnotesize},
        row{4,5,9,13} = {font=\footnotesize},
		cell{6-7}{2-7} = {font=\footnotesize},
        row{2,Y} = {bg=SuCielo},
        colspec = {Q[3,l]Q[1,c]Q[1,c]Q[1,c]Q[1,c]Q[1,c]Q[1,c]},
		row{3,8} = {bg=gray!25},
		cell{4,5,9}{2-7} = {bg=gray!10},
		cell{6-7,10-11}{1} = {bg=gray!10},
    }
    	\SetCell[c=7]{l} ESTADO DEL SENSOR INTERNO \\
    	\SetCell[c=7]{l} ACELERÓGRAFO EN TERRENO NATURAL \\
    	\SetCell[c=7]{l} CENTRADO DE MASA EN LOS TRES CANALES \\
    	\SetCell[r=2]{} &
		\SetCell[c=3]{c} ANTES DEL MANTENIMIENTO &&&
		\SetCell[c=3]{c} DESPUES DEL MANTENIMIENTO \\
    	& HNZ & HNN & HNE & HNZ & HNN & HNE\\
    	NIVEL DE OFFSET &
		\calcularSemisuma{\AntesOffZMax}{\AntesOffZMin} &
		\calcularSemisuma{\AntesOffNMax}{\AntesOffNMin} &
		\calcularSemisuma{\AntesOffEMax}{\AntesOffEMin} &
		\calcularSemisuma{\DespuesOffZMax}{\DespuesOffZMin} &
		\calcularSemisuma{\DespuesOffNMax}{\DespuesOffNMin} &
		\calcularSemisuma{\DespuesOffEMax}{\DespuesOffEMin} 
		\\
    	VALOR EFICAZ - RMS &
		\AntesEficazZ &
		\AntesEficazN &
		\AntesEficazE &
		\DespuesEficazZ &
		\DespuesEficazN &
		\DespuesEficazE
		\\
    	\SetCell[c=7]{l} CALIBRACIÓN DEL SENSOR \\
    	& \SetCell[c=2]{c} FORMA DE ONDA &&
		\SetCell[c=2]{c} AMPLITUD &&
		\SetCell[c=2]{c} FRECUENCIA \\
    	SEÑAL DE CALIBRACIÓN &
		\SetCell[c=2]{c} Sinusoidal &&
		\SetCell[c=2]{c} 1.0 $V_{DC}$ &&
		\SetCell[c=2]{c} 0.25 Hz \\
    	SEÑAL CORREGIDA PROMEDIO &
		\SetCell[c=2]{c} Sinusoidal &&
		\SetCell[c=2]{c} \qty{\AmplitudCorregidaPromedio}{\centi\metre\per\square\second} &&
		\SetCell[c=2]{c} 0.25 Hz \\
    	\SetCell[c=7]{l} OBSERVACIONES\\
    	\SetCell[c=7]{l, 17.5cm}
    	\begin{enumerate}
			\item No corrigió el nivel offset en los tres canales, ya que el valor promedio en cada canal se encontraba muy por debajo del umbral de disparo [1\%g = 10mg]
			\item Con el archivo de calibración se procedió a realizar la corrección instrumental y verificar su valor en cuentas, siendo los valores en cada eje los siguientes:
			\begin{itemize}
				\item VERTICAL: [\num{\CorreccionVerticalMin}; \num{\CorreccionVerticalMax}] - PGA: \qty{\CorreccionVerticalPga}{\centi\metre\per\square\second}
				\item NORTE-SUR: [\num{\CorreccionNorteMin}; \num{\CorreccionNorteMax}] - PGA: \qty{\CorreccionNortePga}{\centi\metre\per\square\second}
				\item ESTE-OESTE: [\num{\CorreccionEsteMin}; \num{\CorreccionEsteMax}] - PGA: \qty{\CorreccionEstePga}{\centi\metre\per\square\second}
			\end{itemize}
    		\item Con el espectrograma se obtiene la frecuencia de la onda sinusoidal registrada en 0.25Hz.
    		\item Para la reconstrución de la señal en términos de aceleración se utilizó la respuesta instrumental mostrada en el <Apartado 11: Formato de Polos y Ceros> de este informe.
    	\end{enumerate}
    \end{tblr}
\end{table}

% \begin{table}[H]
%     \begin{tblr}{
%         cell{6-7,10-12,14-15}{1} = {font=\footnotesize},
%         row{4,5,9,13} = {font=\footnotesize},
%         row{2} = {bg=SuAzul!50},
%         row{3,8,16} = {bg=gray!25},
%         cell{6-7,10-12,14-15}{1} = {bg=gray!25},
%         cell{4,5,9,13}{2-7} = {bg=gray!25},
%         colspec = {Q[3,l]Q[1,c]Q[1,c]Q[1,c]Q[1,c]Q[1,c]Q[1,c]},
%         cell{6-7}{2-7} = {font=\footnotesize},
%     }
%     	\SetCell[c=7]{l} ESTADO DEL SENSOR INTERNO\\
%     	\SetCell[c=7]{l} ACELERÓGRAFO AZOTEA\\
%     	\SetCell[c=7]{l} CENTRADO DE MASA EN LOS TRES CANALES\\
%     	\SetCell[r=2]{} & \SetCell[c=3]{c} ANTES DEL MANTENIMIENTO &&& \SetCell[c=3]{c} DESPUES DEL MANTENIMIENTO\\
%     	& HNZ & HNN & HNE & HNZ & HNN & HNE\\
%     	NIVEL DE OFFSET (count) &
% 		\calcularSemisuma{\oaazmax}{\oaazmin} & \calcularSemisuma{\oaanmax}{\oaanmin} & \calcularSemisuma{\oaaemax}{\oaaemin} & \calcularSemisuma{\oadzmax}{\oadzmin} & \calcularSemisuma{\oadnmax}{\oadnmin} & \calcularSemisuma{\oademax}{\oademin}\\
%     	DIFERENCIA PICO A PICO &
% 		\oaazdif & \oaandif & \oaaedif & \oadzdif & \oadndif & \oadedif\\
%     	\SetCell[c=7]{l} CALIBRACIÓN DEL SENSOR\\
%     	& \SetCell[c=2]{c} FORMA DE ONDA && \SetCell[c=2]{c} AMPLITUD && \SetCell[c=2]{c} ANCHO DE PULSO\\
%     	SEÑAL DE CALIBRACIÓN INYECTADA (EJE Z) & \SetCell[c=2]{c} Pulso && \SetCell[c=2]{c} 0.80 g && \SetCell[c=2]{c} 5.0 seg\\
%     	SEÑAL DE CALIBRACIÓN INYECTADA (EJE Y) & \SetCell[c=2]{c} Pulso && \SetCell[c=2]{c} 0.80 g && \SetCell[c=2]{c} 5.0 seg\\
%     	SEÑAL DE CALIBRACIÓN INYECTADA (EJE X) & \SetCell[c=2]{c} Pulso && \SetCell[c=2]{c} 0.80 g && \SetCell[c=2]{c} 5.0 seg\\
%     	& \SetCell[c=2]{c} CANAL VERTICAL && \SetCell[c=2]{c} CANAL NORTE-SUR && \SetCell[c=2]{c} CANAL ESTE-OESTE\\
%     	SEÑAL RECONSTRUIDA EN CUENTAS (count) &
% 		\SetCell[c=2]{c} \num{1687945.29} &&
% 		\SetCell[c=2]{c} \num{1674033.84} &&
% 		\SetCell[c=2]{c} \num{1696944.57} \\
%     	SEÑAL RECONSTRUIDA EN ACELERACIÓN (g) &
% 		\SetCell[c=2]{c} \num{0.786423} &&
% 		\SetCell[c=2]{c} \num{0.793545} &&
% 		\SetCell[c=2]{c} \num{0.804486} \\
%     	\SetCell[c=7]{l} OBSERVACIONES\\
%     	\SetCell[c=7]{l, 17.5cm}
%     	\begin{enumerate}
%     		\item El centrado de masa se realizó aplicando un filtro pasa alto IIR que establece la frecuencia de corte a 0.01Hz.
%     		\item Para la reconstrución de la señal en términos de aceleración se utilizó la respuesta instrumental mostrada en el <Apartado 11: Formato de Polos y Ceros> a este informe.
%     	\end{enumerate}
%     \end{tblr}
% \end{table}


\begin{longtable}{|c|p{6cm}|c|p{2cm}|*{6}{c|}}
\hline

\rowcolor{SuAzul}
\multicolumn{10}{|l|}{{\color{white} ESTADO DEL SENSOR INTERNO}} \\ \hline

\rowcolor{SuCielo}
\multicolumn{10}{|l|}{VALIDACIÓN CON SISMOS REGISTRADOS} \\ \hline

\rowcolor{gray!25}
\multirow[t]{2}{*}{\footnotesize Ev.} &
\multirow[t]{2}{*}{\footnotesize REFERENCIA} &
\multirow[t]{2}{*}{\footnotesize Mg.} &
\multirow[t]{2}{*}{\footnotesize HORA} &
\multicolumn{3}{c|}{\footnotesize PGA SÓTANO (mm/s$^2$)} &
\multicolumn{3}{c|}{\footnotesize PGA AZOTEA (mm/s$^2$)} \\ \cline{5-7}

\rowcolor{gray!25} &&&&
E-W & N-S & V &
E-W & N-S & V \\ \hline

% \csvreader[
%     head to column names,
%     separator=semicolon,
%     late after line=\\\hline
% ]{ondas/pga.csv}{}{
%     \evento &
%     \referencia &
%     \magnitud &
%     \hora &
%     \sote &q
%     \sotn &
%     \sotz &
%     \azoe &
%     \azon &
%     \azoz
% }

\rowcolor{SuCielo}
\multicolumn{10}{|l|}{OBSERVACIONES} \\ \hline

\multicolumn{10}{|p{17.5cm}|}{\vspace{0.25em}
\begin{minipage}{\linewidth}
	\begin{itemize}
		\item Se presentan los sismos publicados por el Instituto Geofísico del Perú - IGP durante el período de descarga [\DiaIni/\MesIni/\AnioIni~al \DiaMantto/\MesMantto/\AnioMantto] Fuente: [{\color{SuAzul}\url{https://ultimosismo.igp.gob.pe/productos/reportes-sismicos}}]
		\item Se observa que el sismo más significativo dentro del período de descarga ha sido el acontecido el 26/02/2026 18:21:08 a 36km al O de Chilca, Cañete - Lima con una magnitud de 5.0 ML habiendo generado una aceleración máxima de 551 mm/s2 en la componente Vertical.
		\item Las formas de onda de los acelerogramas se presentan en el Apartado 11 <Visualización de registros sísmicos> donde las unidades de aceleración están dados en [mm/s2] para una mejor visualización.
	\end{itemize}
\end{minipage}
} \\ \hline

\end{longtable}